% !TeX root = 03_intro_ml_polytech.tex
% !TeX encoding = UTF-8
% !TeX spellcheck = en_GB
\documentclass[t,10pt,aspectratio=169,rounded]{beamer}
\usetheme{CEA23}

\usepackage[utf8x]{inputenc}
\usepackage[english]{babel}
\usepackage{csquotes}
\usepackage{physics}
\usepackage{svg}

% Tikz styles
\usetikzlibrary{shadows, fit, backgrounds, positioning}
\tikzset{
  % Base style: rounded rectangle with a soft drop shadow
  set/.style = {
    rectangle,
    rounded corners = 8pt,
    draw = CEAblue,
    line width = 0.8pt,
    fill = CEAblue!5!white,
    inner sep = 12pt,
    drop shadow = {shadow xshift = 2pt, shadow yshift = -2pt,
                   opacity = 0.35, fill = black},
  },
  % Colour variants (inherit everything from `set`)
  set green/.style = {set, fill = CEAgreen!5!white, draw = CEAgreen},
  set red/.style   = {set, fill = CEAred!5!white,   draw = CEAred},
  % Label anchored at the upper-left corner, inside the set
  set label/.style = {
    anchor = south west,
    inner sep = 4pt,
    font = \bfseries\small,
  },
  % Element dot
  elem/.style = {circle, fill = black, inner sep = 1.2pt},
}

% Code listings: Carbon-style "code snapshot" cards (see code_listings_demo.tex)
\usepackage{beamer-code-listings}

\AtBeginSection[]{\sepframe[\insertsectionnumber]{}{}}

\begin{document}

\title{Introduction to Machine Learning}
\subtitle{Module 3 -- The Algorithm Toolbox \\[0.5em] SVM, Trees, PCA, and Clustering}
\date[07/12/2026]{7 December 2026}
\author{Riccardo \textsc{Finotello}}

\titlepagelogo{\includegraphics[height=24mm]{graphics/ORIGINAL/ISAS.png}}
\titlepagepicture{img/titlepage.png}
\setvalue{\diffusion}{}

\seminar[PolytechParis]{Machine Learning - PolytechParis}

\CEAmerci{See you next time!}
\CEAcenter{CEA SACLAY}
\addr{91\,191 Gif-sur-Yvette Cedex\\France}
\CEAemail{\href{mailto:riccardo.finotello@cea.fr}{\texttt{riccardo.finotello@cea.fr}}}

% --- TITLE SLIDE ---
\begin{frame}[plain]
	\titlepage
\end{frame}

% ------------------------------------------------------------------
% Content slides: drop your module slide fragments into 03_slides/
% and include them below, e.g.
%   \input{03_slides/section_01}
%   \input{03_slides/section_02}
% Each fragment contains \section{...} and \begin{frame}...\end{frame}
% only (no preamble, no \begin{document}).
% ------------------------------------------------------------------

% --- CONCLUDING SLIDE ---
\begin{frame}[plain]
	\usebeamertemplate*{last page}
\end{frame}
\end{document}
